\documentclass[10pt, a4paper]{article}

% Parametri che modificano il file main.tex
% Le uniche parti da cambiare su main.tex sono:
% - vari \vspace tra sezioni
% - tabella azioni da intraprendere
% - sezione altro

\def\data{2023-11-09}
\def\oraInizio{17:00}
\def\oraFine{18:30}
\def\luogo{Sede di AzzurroDigitale: \\ Via della Croce Rossa, 42 (PD)}

\def\tipoVerb{Esterno} % Interno - Esterno

\def\nomeResp{Feltrin E.} % Cognome N.
\def\nomeVer{Campese M.} % Cognome N.
\def\nomeSegr{Ciriolo I.} % Cognome N.

\def\nomeAzienda{Azzurro Digitale}
\def\firmaAzienda{azzurrodigitale.png}
\def\firmaResp{emanuele.png} % nome Responsabile

\def\listaPartInt{
Bresolin G.,
Campese M.,
Ciriolo I.,
Dugo A.,
Feltrin E.,
Michelon R.,
Orlandi G.
}

\def\listaPartEst{
\textit{Azzurro Digitale:},
Caliendo G.,
D'Avanzo C.,
Picello M., \\
\textit{Coding Cowboys:},
Cecchin A.,
Dolzan M.,
Ferraioli F.,
Giacomuzzo F.,
Lago L.,
Menon G.,
Nordio A.
}

\def\listaRevisioneAzioni {x}

\def\listaOrdineGiorno {
{Visita della sede dell'azienda AzzurroDigitale;},
{Breve presentazione dell'azienda da parte di Picello C., specialista in People\&Culture, con introduzione ai valori aziendali di AzzurroDigitale;},
{Possibilità di porre alcune domande generiche sul progetto ai referenti Caliendo G. e D'Avanzo C., con seguente approfondimento dei vari temi toccati.\\ All'incontro ed alle discussioni hanno partecipato anche i membri del gruppo \textit{Coding Cowboys}.}
}

\def\listaDiscussioneInterna {
{Discussione 1;},
{Discussione 2;},
}

\newcommand{\domris}[2]{\textbf{#1}\\#2}

\def\listaDiscussioneEsterna {
\domris
{Quali riferimenti dei membri del gruppo è preferibile inoltrare all'azienda?
}
{Sono necessari gli account che devono avere l'accesso al \textit{repository} \textit{GitHub} fornito dall'azienda, un indirizzo e-mail del team per l'accesso al canale \textit{Slack} fornito dall'azienda ed i singoli indirizzi e-mail dei componenti;
},
\domris
{Quando sarà possibile discutere nello specifico dell'analisi dei requisiti del progetto?
}
{L'azienda aveva già in programma un incontro in data 2023-11-15 della durata di 1.30h per la discussione dei requisiti del progetto con ciascun gruppo preso singolarmente. I gruppi hanno garantito entrambi disponibilità a tali incontri;
},
\domris
{L'azienda in precedenza ha consigliato l'adozione di un framework di tipo \textit{agile} (nello specifico di \textit{Scrum}) per lo sviluppo del progetto: dal momento che i membri del team non hanno mai avuto la possibilità di sperimentarlo, l'azienda propone qualche consiglio per la sua applicazione?
}
{L'azienda ha fornito una spiegazione dei vari ruoli all'interno di un framework \textit{Scrum} generico fornendo al team le basi per iniziare ad apprenderlo ed applicarlo, specificando che si tratta una metodologia complessa che si apprende con maggior successo durante la sua applicazione. Viene sottolineata l'importanza di avere un solido sistema di \textit{backlog} per l'accertamento dei progressi del progetto;
},
\domris
{Quali criteri utilizzerà l'azienda per la valutazione del successo del progetto?
}
{È stato spiegato che anche questo aspetto del progetto sarà analizzato in futuro con il team stesso, poiché al momento non vi è ancora un'idea precisa;
},
\domris
{Cosa intende di preciso l'azienda riferendosi ad un sistema di \textit{Bug Reporting} (e alla sua relativa documentazione)?
}
{L'azienda ha chiarito che la documentazione non è l'interesse principale e che come sistema di \textit{Bug Reporting} è sufficiente servirsi di un\textit{ Issue Tracking System}.
}
} 




\def\listaDecisioni {x}

\usepackage{style}
\usepackage{headerfooter}

\title{\titolo}
\author{SWEetCode}

\begin{document}

% PRIMA PAGINA
\begin{titlepage}
    \thispagestyle{empty}
    \begin{tikzpicture}[remember picture, overlay]
        % TRIANGOLI
        \draw[fill=secondarycolor, secondarycolor] (current page.north west) -- (current page.south west) -- (8.8, -28);
        \draw[fill=primarycolor, primarycolor] (-3, 5) -- (4, -13.6) -- (11, 5);

        % LOGO
        \node [xshift=-5cm, yshift=25cm] (logo) at (current page.south east) {\includegraphics[width=6.5cm]{img/logo.png}};

        % SWEETCODE - DATE
        \node [anchor=north east, align=right, xshift=-1.2cm, yshift=20.5cm, text=black] (sweetcode) at (current page.south east) {\fontsize{32pt}{36pt}\selectfont SWEetCode};
        \draw[line width=4pt, lightcol] ([xshift=-3cm, yshift=-0.37cm]sweetcode.south west) -- ([yshift=-0.37cm]sweetcode.south east);
        \node [anchor=north east, align=right, xshift=-1.2cm, yshift=18.7cm, text=black] (date) at (current page.south east){\fontsize{24pt}{24pt} \selectfont Verbale \tipoVerb};

        % NOME FILE
        \node [anchor=north east, text width=15cm, align=right, xshift=-1.2cm, yshift=17cm, text=black] (titolo) at (current page.south east){\fontsize{48pt}{48pt}\textbf{\data}};

        % BOX DATI PARTECIPANTI
        \ifthenelse{\equal{\tipoVerb}{Esterno}}{
            \node[anchor=north east, xshift=-1.2cm, yshift=14.5cm, minimum width=8cm] (box) at (current page.south east){};
        }{
            \node[anchor=north east, xshift=-1.2cm, yshift=12.5cm, minimum width=8cm] (box) at (current page.south east){};
        }

        % RESPONSABILE
        \node[anchor=north west, align=left] (dati2) at (box.north west) {\fontsize{15pt}{15pt}\selectfont \textbf{Responsabile}};
        \draw[line width=4pt, lightcol] (dati2.south west) -- ([xshift=8cm]dati2.south west);
        \node[anchor=north west, align=left] (dati21) at (dati2.south west){\fontsize{13pt}{13pt}\selectfont \nomeResp};

        % VERIFICATORE
        \node[anchor=north west, yshift=-1cm, align=left] (dati3) at (dati21.north west) {\fontsize{15pt}{15pt}\selectfont \textbf{Verificatore}};
        \draw[line width=4pt, lightcol] (dati3.south west) -- ([xshift=8cm]dati3.south west);
        \node[anchor=north west, align=left] (dati31) at (dati3.south west){\fontsize{13pt}{13pt}\selectfont \nomeVer};

        % SEGRETARIO DI RIUNIONE
        \node[anchor=north west, yshift=-1cm, align=left] (dati4) at (dati31.north west) {\fontsize{15pt}{15pt}\selectfont \textbf{Segretario di Riunione}};
        \draw[line width=4pt, lightcol] (dati4.south west) -- ([xshift=8cm]dati4.south west);
        \node[anchor=north west, align=left] (dati41) at (dati4.south west){\fontsize{13pt}{13pt}\selectfont \nomeSegr};
        
        % UNIPD - SWE
        \node [xshift=4.4cm, yshift=2.3cm, draw, secondarycolor, text=white] (uni) at (current page.south west) {\fontsize{20pt}{20pt} \selectfont Università di Padova};
        \node [xshift=0.65cm, yshift=0.7cm, draw, secondarycolor, text=white, below=of uni] (corso) {\fontsize{20pt}{20pt}\selectfont Ingegneria del Software};

        % FIRMA AZIENDA
        \ifthenelse{\equal{\tipoVerb}{Esterno}}{
            \draw[line width=4pt, lightcol] ([xshift=-1.2cm, yshift=7cm]current page.south east) -- ([xshift=-8cm, yshift=7cm]current page.south east);
            \node [xshift=-4.8cm, yshift=8cm] (logo) at (current page.south east) {\includegraphics[width=6cm]{img/firme/\firmaAzienda}};
            \node[anchor=north west, xshift=12.9cm, yshift=6.65cm, align=left] at (current page.south west)
            {\fontsize{13pt}{13pt}\selectfont L'azienda: \nomeAzienda};
        }

        % FIRMA CODING COWBOYS
        \draw[line width=4pt, lightcol] ([xshift=-1.2cm, yshift=4.35cm]current page.south east) -- ([xshift=-8cm, yshift=4.35cm]current page.south east);
        \node [xshift=-4.8cm, yshift=5.5cm] (logo) at (current page.south east) {\includegraphics[width=6cm]{img/firme/menon.png}};
        \node[anchor=north west, xshift=12.9cm, yshift=4cm, align=left] at (current page.south west)
        {\fontsize{13pt}{13pt}\selectfont Coding Cowboys: Menon G.};

        % FIRMA
        \draw[line width=4pt, lightcol] ([xshift=-1.2cm, yshift=1.8cm]current page.south east) -- ([xshift=-8cm, yshift=1.8cm]current page.south east);
        \node [xshift=-4.8cm, yshift=2.45cm] (logo) at (current page.south east) {\includegraphics[width=6cm]{img/firme/\firmaResp}};
        \node[anchor=north west, xshift=12.9cm, yshift=1.45cm, align=left] at (current page.south west)
        {\fontsize{13pt}{13pt}\selectfont Il Responsabile: \nomeResp};
        
    \end{tikzpicture}
\end{titlepage}

% REGISTRO DELLE VERSIONI
{\renewcommand{\arraystretch}{1.5}
\section*{Registro delle versioni}
\begin{tabularx}{\textwidth}{c|c|c|c|X}
\textbf{Versione} & \textbf{Data} & \quantities{\textbf{Responsabile di}\\\textbf{stesura}}& \textbf{Revisore} & \quantities{\textbf{Dettaglio e}\\\textbf{motivazioni}} \\
\hline
v1.0.0(3) & $2023-11-17$ & Michelon R. & Bresolin G. & Correzione paragrafo scopo del documento.\\
\hline
v0.0.1(22) & $2023-11-13$ & Bresolin G. & Campese M. & Aggiornamento struttura di documento: modifica ordine registro delle versioni dalla più recente e aggiunta voce ”motivazioni”; verbali: introduzione \#issue nella tabella delle azioni.\\
\hline
v0.0.1(21) & $2023-11-13$ & Bresolin G. & Campese M. & Aggiornamento documentazione: struttura di documento: presenza obbligatoria registro delle versioni e rimozione versione;  norme tipografiche: aggiornamento versioni. Aggiunta configuration managment: version control.\\
\hline
v0.0.1(20) & $2023-10-30$ & Bresolin G. & Michelon R. & Processi di supporto: configuration management: \textit{ITS} e pull request; Processi organizzativi: \textit{GitHub Teams}.\\ 
\hline
v0.0.1(19) & $2023-10-30$ & Campese M. & Michelon R. & Documentazione: strutture di documento: posta elettronica.\\
\hline
v0.0.1(18) & $2023-10-30$ & Dugo A. & Bresolin G. & Processi organizzativi: comunicazione.\\
\hline
v0.0.1(13) & $2023-10-28$ & Ciriolo I. & Campese M. & Strumenti.\\
\hline
v0.0.1(8) & $2023-10-27$ & Bresolin G. & Ciriolo I. & Processo di supporto: documentazione e verifica.\\
\hline
v0.0.1(5) & $2023-10-26$ & Campese M. & Orlandi G. &  Introduzione.\\

\end{tabularx}}

\newpage

% INDICE
\tableofcontents
\newpage

% INTRODUZIONE
\section{Introduzione}
\subsection{Scopo del documento}
Il Piano di Qualifica è un documento che si valuta di modificare incrementalmente, in particolare per 
definire le metriche di valutazione del prodotto, che saranno definite conformemente ai requisiti e 
alle aspettative del proponente, al fine di poter correttamente definire la qualità del prodotto, 
attraverso un processo di miglioramento continuo e che, per sua natura, tende a diventare 
incrementale nel corso del tempo e quando viene definita baseline. Per tutti questi motivi, la qualità 
viene definita da un insieme di processi che cerchino di definire metriche di misurazione di efficacia 
ed efficienza (misure quantitative che serviranno da valutazione nel corso di realizzazione del 
progetto didattico)
\paragraph{}
%\subsection{Scopo del prodotto}
\subsection{Glossario}
Al fine di evitare incomprensioni relative alla terminologia usata all'interno del documento, viene 
fornito un Glossario nel file apposito, tale da non avere terminologie ambigue nell'attività 
progettuale individuata e dandone una definizione precisa. Ogni termine avrà nel documento una 
lettera G come apice, per meglio evidenziare la loro appartenenza al documento indicato.

\subsection{Riferimenti}
   \subsubsection{Riferimenti normativi}
   \begin{itemize}
    \item \textit{Regolamento del progetto didattico}: \\
    \href{https://www.math.unipd.it/~tullio/IS-1/2023/Dispense/PD2.pdf}{https://www.math.unipd.it/~tullio/IS-1/2023/Dispense/PD2.pdf}\\
    \item \textit{Norme di progetto v0.4.3};
    \end{itemize}
    
    \subsubsection{Riferimenti informativi}
    
    \begin{itemize}
    \item \textit{Glossario v0.0.1} (da creare parallelamente) 
    \item \textit{Presentazione capitolato C1 - AzzurroDigitale}:\\
    \href{https://www.math.unipd.it/~tullio/IS-1/2023/Progetto/C1.pdf}{https://www.math.unipd.it/~tullio/IS-1/2023/Progetto/C1.pdf}
    \item \textit{Verbali esterni ed interni};
    \item \textit{Analisi e descrizione delle funzionalità: Use Case e relativi diagrammi (UML)}:\\
    \href{https://www.math.unipd.it/~rcardin/swea/2022/Diagrammi%20Use%20Case.pdf}{https://www.math.unipd.it/~rcardin/swea/2022/Diagrammi\%20Use\%20Case.pdf}
    \item \textit{Analisi e descrizione delle funzionalità: Diagrammi delle Attività (UML)}:\\
    \href{https://www.math.unipd.it/~rcardin/swea/2022/Diagrammi%20di%20Attivit%C3%A0.pdf}{https://www.math.unipd.it/~rcardin/swea/2022/Diagrammi\%20di\%20Attivit\%C3\%A0.pdf}
    \item \textit{Progettazione e Programmazione: Diagrammi delle Classi}:\\
    \href{https://www.math.unipd.it/~rcardin/swea/2023/Diagrammi%20delle%20Classi.pdf}{https://www.math.unipd.it/~rcardin/swea/2023/Diagrammi\%20delle\%20Classi.pdf}
    \end{itemize}

% Qualità di Processo
\newpage
\section{Qualità di processo}
\subsection{Scopo ed obiettivi}
La qualità è determinata univocamente dai processi che lo compongono, misurata mettendo in atto 
delle metriche che permettano di valutare tali processi e accertarsi che raggiungano i corretti 
obiettivi di qualità previsti. In particolare, si fa riferimento al cosiddetto Ciclo PDCA (Plan - Do - Check 
- Act), atto a garantire un miglioramento continuo nell'utilizzo dei processi e delle risorse tramite 
pianificazione, successiva verifica con le metriche previste ed integrazione in fase di produzione. Di 
seguito, i processi individuati e i livelli di qualità previsti per ciascuno.
In particolare, per ciascuna metrica si opera una breve descrizione, dando un'idea comprensiva 
dell'attuazione e dei valori considerati accettabili in fase di controllo (check) qualità

\subsection{Processi primari}
    {\renewcommand{\arraystretch}{1.5}
    \begin{tabularx}{\textwidth}{p{0.18\textwidth}|p{0.6\textwidth}|X}
    \textbf{Obiettivo} & \textbf{Descrizione} & \textbf{Metriche}  \\
    \hline
    Fornitura &  & \\
    \hline
    Sviluppo &  &  \\
    \end{tabularx}}

\subsection{Processi di supporto}
    {\renewcommand{\arraystretch}{1.5}
    \begin{tabularx}{\textwidth}{p{0.18\textwidth}|p{0.6\textwidth}|X}
    \textbf{Obiettivo} & \textbf{Descrizione} & \textbf{Metriche}  \\
    \hline
    Verifica &  &  \\
    \hline
    Gestione della qualità &  &  \\
    \end{tabularx}}
    
\subsection{Processi organizzativi}
    {\renewcommand{\arraystretch}{1.5}
    \begin{tabularx}{\textwidth}{p{0.18\textwidth}|p{0.6\textwidth}|X}
    \textbf{Obiettivo} & \textbf{Descrizione} & \textbf{Metriche}  \\
    \hline
    Gestione organizzativa &  &  \\
    \end{tabularx}}
    
\subsection{Metriche utilizzate}
\subsubsection{Processi primari}
    {\renewcommand{\arraystretch}{1.5}
    \begin{tabularx}{\textwidth}{p{0.05\textwidth}|p{0.35\textwidth}|X|X}
    \textbf{ID} & \textbf{Nome metrica} & \textbf{Valore accettabile} & \textbf{Valore ottimale}  \\
    \hline
    \multicolumn{4}{l}{\cellcolor{primarycolor}\textbf{\textit{Fornitura}}} \\
    \hline
     &  &  &  \\
    \hline
     &  &  &  \\
    \hline
     &  &  &  \\
    \hline
     &  &  &  \\
    \hline
     &  &  &  \\
    \hline
     &  &  &  \\
    \hline
     &  &  &  \\
    \hline
    \multicolumn{4}{l}{\cellcolor{primarycolor}\textbf{\textit{Sviluppo}}} \\
    \hline
     &  &  &  \\
    \hline
     &  &  &  \\
    \end{tabularx}}
   
\subsubsection{Processi di supporto}
    {\renewcommand{\arraystretch}{1.5}
    \begin{tabularx}{\textwidth}{p{0.05\textwidth}|p{0.35\textwidth}|X|X}
    \textbf{ID} & \textbf{Nome metrica} & \textbf{Valore accettabile} & \textbf{Valore ottimale}  \\
    \hline
    \multicolumn{4}{l}{\cellcolor{primarycolor}\textbf{\textit{Verifica}}} \\
    \hline
     &  &  &  \\
    \hline
     &  &  &  \\
    \hline
     &  &  &  \\
    \hline
     &  &  &  \\
    \hline
    \multicolumn{4}{l}{\cellcolor{primarycolor}\textbf{\textit{Gestione della qualità}}} \\
    \hline
     &  &  &  \\
    \end{tabularx}}    
    
\subsubsection{Processi organizzativi}
    {\renewcommand{\arraystretch}{1.5}
    \begin{tabularx}{\textwidth}{p{0.05\textwidth}|p{0.35\textwidth}|X|X}
    \textbf{ID} & \textbf{Nome metrica} & \textbf{Valore accettabile} & \textbf{Valore ottimale}  \\
    \hline
    \multicolumn{4}{l}{\cellcolor{primarycolor}\textbf{\textit{Gestione organizzativa}}} \\
    \hline
     &  &  &  \\
    \end{tabularx}}  

% Qualità di prodotto
\newpage
\section{Qualità di prodotto}
Dopo un'attenta analisi per individuare le proprietà utili per la gestione del ciclo di vita del software 
si è cercato di trovare quali caratteristiche siano necessarie per la realizzazione di un prodotto di 
qualità.
L'acronimo principale di riferimento è MPD, cioè Mean Percentage Difference (MPD), metrica 
utilizzata per valutare la differenza percentuale media tra due valori. Ad esempio, può essere 
utilizzata per calcolare la differenza percentuale media tra i valori di previsione e i valori effettivi di 
una variabile in un modello di previsione o di apprendimento automatico.
In altre parole, l'MPD indica la percentuale media di errore tra le previsioni del modello e i valori 
effettivi della variabile in questione. Un valore di MPD inferiore indica che le previsioni del modello 
sono più precise, mentre un valore più elevato indica una maggiore imprecisione.
\subsection{Documentazione}
    {\renewcommand{\arraystretch}{1.5}
    \begin{tabularx}{\textwidth}{p{0.18\textwidth}|p{0.6\textwidth}|X}
    \textbf{Obiettivo} & \textbf{Descrizione} & \textbf{Metriche}  \\
    \hline
    Correttezza linguistica &  & \\
    \hline
    Leggibilità &  &  \\
    \end{tabularx}}
    
\subsection{Software}
    {\renewcommand{\arraystretch}{1.5}
    \begin{tabularx}{\textwidth}{p{0.18\textwidth}|p{0.6\textwidth}|X}
    \textbf{Obiettivo} & \textbf{Descrizione} & \textbf{Metriche}  \\
    \hline
    Funzionalità &  & \\
    \hline
    Usabilità &  & \\
    \hline
    Efficienza &  &  \\
    \hline
    Affidabilità &  &  \\
    \hline
    Portabilità &  &  \\
    \hline
    Manutenibilità &  &  \\
    \hline
    Copertura dei test &  &  \\
    \end{tabularx}}

\subsection{Metriche utilizzate}
\subsubsection{Documentazione}
    {\renewcommand{\arraystretch}{1.5}
    \begin{tabularx}{\textwidth}{p{0.05\textwidth}|p{0.35\textwidth}|X|X}
    \textbf{ID} & \textbf{Nome metrica} & \textbf{Valore accettabile} & \textbf{Valore ottimale}  \\
    \hline
    \multicolumn{4}{l}{\cellcolor{primarycolor}\textbf{\textit{Correttezza linguistica}}} \\
    \hline
     &  &  &  \\
    \hline
    \multicolumn{4}{l}{\cellcolor{primarycolor}\textbf{\textit{Leggibilità}}} \\
    \hline
     &  &  &  \\
    \end{tabularx}}
    
\subsubsection{Software}
    {\renewcommand{\arraystretch}{1.5}
    \begin{tabularx}{\textwidth}{p{0.05\textwidth}|p{0.35\textwidth}|X|X}
    \textbf{ID} & \textbf{Nome metrica} & \textbf{Valore accettabile} & \textbf{Valore ottimale}  \\
    \hline
    \multicolumn{4}{l}{\cellcolor{primarycolor}\textbf{\textit{Funzionalità}}} \\
    \hline
     &  &  &  \\
    \hline
     &  &  &  \\
    \hline
     &  &  &  \\
    \hline
    \multicolumn{4}{l}{\cellcolor{primarycolor}\textbf{\textit{Usabilità}}} \\
    \hline
     &  &  &  \\
    \hline
     &  &  &  \\
    \multicolumn{4}{l}{\cellcolor{primarycolor}\textbf{\textit{Efficienza}}} \\
    \hline
     &  &  &  \\
    \hline
    \multicolumn{4}{l}{\cellcolor{primarycolor}\textbf{\textit{Affidabilità}}} \\
    \hline
     &  &  &  \\
    \hline
    \multicolumn{4}{l}{\cellcolor{primarycolor}\textbf{\textit{Portabilità}}} \\
    \hline
     &  &  &  \\
    \hline
    \multicolumn{4}{l}{\cellcolor{primarycolor}\textbf{\textit{Manutenibilità}}} \\
    \hline
     &  &  &  \\
    \hline
    \multicolumn{4}{l}{\cellcolor{primarycolor}\textbf{\textit{Copertura dei test}}} \\
    \hline
     &  &  &  \\
    \end{tabularx}}

%Test e specifiche
\newpage
\section{Test e specifiche}
Nella seguente sezione verranno espresse in maniera dettagliata le varie metodologie di test, gli 
obiettivi del testing e i criteri di successo utilizzati durante lo sviluppo del prodotto.
Il gruppo SWEetCode per perseguire la correttezza del prodotto e facilitare la fase di validazione, svolgerà la verifica in parallelo allo sviluppo (Modello a VG).
I test dovranno essere resi il più automatici possibile, per evitare che la fase di testing rallenti la produzione.\\
Lo stato di ciascun test può avere due valori:
\begin{itemize}
\item I: Implementato;
\item NI: Non Implementato.
\end{itemize}

\subsection{Formato ID dei test}
\subsection{Test di unità}
    {\renewcommand{\arraystretch}{1.5}
    \begin{tabularx}{\textwidth}{p{0.12\textwidth}|p{0.7\textwidth}|X}
    \textbf{ID} & \textbf{Descrizione} & \textbf{Stato}  \\
    \hline
     &  & \\
    \hline
     &  &  \\
    \hline
     &  & \\
    \hline
     &  &  \\
    \end{tabularx}}
    
\subsection{Test di integrazione}
    {\renewcommand{\arraystretch}{1.5}
    \begin{tabularx}{\textwidth}{p{0.12\textwidth}|p{0.7\textwidth}|X}
    \textbf{ID} & \textbf{Descrizione} & \textbf{Stato}  \\
    \hline
     &  & \\
    \hline
     &  &  \\
    \hline
     &  & \\
    \end{tabularx}}

\subsection{Test di sistema}
    {\renewcommand{\arraystretch}{1.5}
    \begin{tabularx}{\textwidth}{p{0.12\textwidth}|p{0.55\textwidth}|p{0.12\textwidth}|X}
    \textbf{ID} & \textbf{Descrizione} & \textbf{Requisito} & \textbf{Stato}  \\
    \hline
     &  & \\
    \hline
     &  &  \\
    \end{tabularx}}

\subsection{Test di accettazione}
    {\renewcommand{\arraystretch}{1.5}
    \begin{tabularx}{\textwidth}{p{0.12\textwidth}|p{0.7\textwidth}|X}
    \textbf{ID} & \textbf{Descrizione} & \textbf{Stato}  \\
    \hline
     &  & \\
    \hline
     &  &  \\
    \end{tabularx}}
    
\section{Resoconto delle attività di verifica}

\subsection{Verifica della documentazione}
\subsubsection{Errori ortografici}
\subsubsection{Indici di Gulpease}

\subsection{Verifica dei processi}
\subsubsection{Estimated at completion}
\subsubsection{Budget variance e schedule variance}
\subsubsection{Actual cost e estimate to complete}
\subsubsection{Earned value e planned value}
\subsubsection{Requirements stability index e satisfied obligatory requirements}
\subsubsection{Code coverage back-end}
\subsubsection{Code coverage front-end}
\subsubsection{Passed test cases percentage}
\subsubsection{Failed test cases percentage}
\subsubsection{Comprensibilità del codice}

\end{document}
